% 依赖 booktabs、array、tabularx、multirow 和 threeparttable；保留原稿标签。

\begin{table}[htbp]
  \centering
  \small
  \setlength{\tabcolsep}{3.5pt}
  \renewcommand{\arraystretch}{1.25}
  \begin{threeparttable}
    \caption{启动时易感者人数与拐点相对位置的参数敏感性}
    \label{tab:s1:lambda-sensitivity}
    \begin{tabularx}{\textwidth}{@{}>{\raggedright\arraybackslash}p{0.29\textwidth}*{4}{>{\centering\arraybackslash}p{0.10\textwidth}}>{\raggedright\arraybackslash}X@{}}
      \toprule
      \multirow{2}{*}{分析量} & \multicolumn{4}{c}{偏导数符号} & \multirow{2}{*}{依据}\\
      \cmidrule(lr){2-5}
       & $\beta$ & $q_0$ & $c_0$ & $\eta$ & \\
      \midrule
      启动时易感者人数 $S^*$ & $>0$ & $<0$ & $>0$ & $<0$ & 引理~\ref{lem:s1:Sstar-sensitivity}\\
      拐点相对位置 $\lambda$ & $>0$ & 不定 & $>0$ & $<0$ & 命题~\ref{prop:s1:lambda-sensitivity}\\
      \bottomrule
    \end{tabularx}
    \begin{tablenotes}[flushleft]
      \footnotesize
      \item 注：表中各项为分析量关于列示参数的偏导数符号，求导时固定 $N,S_0,I_0$ 及其余模型参数。关于 $S^*$ 的结果要求满足控制触发条件；关于 $\lambda$ 的结果还要求 $S_c<2\bar S<S^*$。不定表示偏导数的符号不能在上述参数范围内统一确定。
    \end{tablenotes}
  \end{threeparttable}
\end{table}

\begin{table}[htbp]
  \centering
  \small
  \setlength{\tabcolsep}{3.5pt}
  \renewcommand{\arraystretch}{1.25}
  \begin{threeparttable}
    \caption{拐点相关量与控制持续时间的参数敏感性}
    \label{tab:s1:inflection-roles}
    \begin{tabularx}{\textwidth}{@{}>{\raggedright\arraybackslash}p{0.29\textwidth}*{4}{>{\centering\arraybackslash}p{0.10\textwidth}}>{\raggedright\arraybackslash}X@{}}
      \toprule
      \multirow{2}{*}{分析量} & \multicolumn{4}{c}{偏导数符号} & \multirow{2}{*}{依据}\\
      \cmidrule(lr){2-5}
       & $\beta$ & $q_0$ & $c_0$ & $\eta$ & \\
      \midrule
      拐点处隔离率 $\qinf$ & $<0$ & $=0$ & $=0$ & $=0$ & 定理~\ref{thm:s1:inflection}\\
      $\qinf-q_0$ & $<0$ & $<0$ & $=0$ & $=0$ & 式~\eqref{eq:s1:inflection-exist}\\
      $S^*-2\bar S$ & $>0$ & $<0$ & $>0$ & $<0$ & 引理~\ref{lem:s1:Sstar-sensitivity}、式~\eqref{eq:s1:barS}\\
      控制持续时间 $\Delta t$ & $>0$ & $<0$ & 不定 & $<0$ & 推论~\ref{cor:s1:duration-sensitivity}\\
      \bottomrule
    \end{tabularx}
    \begin{tablenotes}[flushleft]
      \footnotesize
      \item 注：表中各项为分析量关于列示参数的偏导数符号，求导时固定 $N,S_0,I_0$ 及其余模型参数。结果在控制触发条件下成立。$\qinf-q_0$ 和 $S^*-2\bar S$ 同时为正时存在内部拐点；$=0$ 表示该量不依赖相应参数。$\partial\Delta t/\partial c_0$ 的符号由推论~\ref{cor:s1:duration-sensitivity} 中的判据确定。
    \end{tablenotes}
  \end{threeparttable}
\end{table}
